\DocumentMetadata{
  pdfversion=2.0,
  pdfstandard=ua-2,
  lang=en-US,
  tagging=on,
  tagging-setup={math/setup=mathml-SE,tabsorder=structure}
}
\documentclass[11pt,letterpaper]{article}
\usepackage[margin=0.68in,includefoot]{geometry}
\usepackage{newcomputermodern}
\usepackage{amsmath,amscd,mathtools}
\usepackage{tikz,adjustbox}
\providecommand{\square}{\mdlgwhtsquare}
\usepackage[hidelinks]{hyperref}
\hypersetup{
  pdftitle={Algebra First-Year Exam, January 2025, Part 1},
  pdfauthor={University of Florida Department of Mathematics},
  pdfsubject={Graduate mathematics examination},
  pdfdisplaydoctitle=true,
  bookmarksopen=true
}
\setcounter{secnumdepth}{0}
\setlength{\parindent}{0pt}
\setlength{\parskip}{0.28em}
\setlength{\emergencystretch}{3em}
\setlength{\leftmargini}{2.15em}
\setlength{\leftmarginii}{2.65em}
\setlength{\leftmarginiii}{3em}
\pagestyle{empty}
\raggedbottom
\allowdisplaybreaks
\NewTaggingSocketPlug{block/list/label}{examproblem}{%
  \tagstructbegin{tag=Lbl}\tagstructend%
  \tagstructbegin{tag=\LBody}%
  \tagstructbegin{tag=\ProblemHeadingTag,title-o={\ProblemHeadingText}}%
  \tagmcbegin{tag=\ProblemHeadingTag,actualtext={\ProblemHeadingText}}%
  #2%
  \tagmcend%
  \tagstructend%
}
\newenvironment{examproblems}[2]{%
  \def\ProblemHeadingTag{#2}%
  \AssignTaggingSocketPlug{block/list/label}{examproblem}%
  \begin{enumerate}%
  \setcounter{enumi}{#1}%
  \renewcommand{\labelenumi}{\textbf{\theenumi.}}%
  \setlength{\topsep}{0.35em}%
  \setlength{\partopsep}{0pt}%
  \setlength{\itemsep}{0.48em}%
  \setlength{\parsep}{0.18em}%
}{\end{enumerate}}
\newenvironment{examsubparts}{%
  \AssignTaggingSocketPlug{block/list/label}{default}%
  \begin{enumerate}%
  \setlength{\topsep}{0.18em}%
  \setlength{\partopsep}{0pt}%
  \setlength{\itemsep}{0.16em}%
  \setlength{\parsep}{0.1em}%
}{\end{enumerate}}
\newenvironment{examinstructions}{%
  \par\begingroup\itshape
}{%
  \par\endgroup\vspace{0.18em}
}
\makeatletter
\renewcommand\subsection{\@startsection{subsection}{2}{0pt}%
  {-1.25ex plus -.25ex minus -.1ex}%
  {0.55ex plus .1ex}%
  {\normalfont\large\bfseries}}
\renewcommand{\@maketitle}{%
  \newpage\null\vskip -1em%
  {\footnotesize\bfseries
    \href{https://www.ufl.edu/}{University of Florida}, \href{https://math.ufl.edu/}{Department of Mathematics}\par}%
  \vskip -0.5em%
  \tagstructbegin{tag=H1,title={Algebra First-Year Exam, January 2025, Part 1}}%
  \tagpdfparaOff%
  \tagmcbegin{tag=H1}%
    {\LARGE\bfseries\href{https://gma.math.ufl.edu/exams/algebra-first-year/}{Algebra First-Year Exam}, January 2025, Part 1\par}%
  \tagmcend%
  \tagpdfparaOn%
  \tagstructend%
  \vskip 0.25em%
  \tagmcbegin{artifact}%
    \hrule height 1pt%
  \tagmcend%
  \vskip 1em%
}
\makeatother
\title{Algebra First-Year Exam, January 2025, Part 1}
\author{}
\date{}
\begin{document}
\maketitle
\thispagestyle{empty}
\begin{examinstructions}
Answer 4 questions. You should indicate which 4 problems you wish to be graded. Write your solutions clearly in complete English sentences. You may quote theorems (within reason) as long as you state them clearly.
\end{examinstructions}
\begin{examproblems}{0}{H2}
\def\ProblemHeadingText{Problem 1.}
\item
In \(S_8\), write down a representative of each conjugacy class of elements of order~\(4\) and compute the order of its centralizer.
\def\ProblemHeadingText{Problem 2.}
\item
Let~\(G\) be a group of order \(124=2^2\cdot 31\). Prove that~\(G\) has a cyclic subgroup of order \(62\).
\def\ProblemHeadingText{Problem 3.}
\item
Prove that there exists a nonabelian group of order \(2025=3^4\cdot 5^2\). (You do not have to determine its multiplication explicitly.)
\def\ProblemHeadingText{Problem 4.}
\item
This problem has three parts.
\begin{examsubparts}
\item[\textnormal{(a)}]
Give the definition of a maximal subgroup of a nontrivial group.
\item[\textnormal{(b)}]
Find two maximal subgroups of \(S_5\) of different orders.
\item[\textnormal{(c)}]
Show that the group \((\mathbb{Q},+)\) has no maximal subgroup.
\end{examsubparts}
\def\ProblemHeadingText{Problem 5.}
\item
Let~\(G\) be a finite group and~\(p\) a prime. Let~\(K\) be the intersection of all the Sylow~\(p\)-subgroups of~\(G\).
\begin{examsubparts}
\item[\textnormal{(a)}]
Prove that~\(K\) is a characteristic subgroup of~\(G\).
\item[\textnormal{(b)}]
Suppose~\(H\) is a normal~\(p\)-subgroup of~\(G\). Prove that \(H\subseteq K\).
\end{examsubparts}
\def\ProblemHeadingText{Problem 6.}
\item
Let~\(G\) be a group, and let \(Z(G)\) denote its center.
\begin{examsubparts}
\item[\textnormal{(a)}]
State what it means for~\(G\) to be nilpotent of class~\(c\).
\item[\textnormal{(b)}]
Prove that if~\(G\) is nilpotent of class~\(1\), then~\(G\) is abelian.
\item[\textnormal{(c)}]
Define the commutator subgroup \(G'\) of~\(G\).
\item[\textnormal{(d)}]
Prove that if~\(G\) is nilpotent of class~\(2\), then \(G'\subseteq Z(G)\).
\end{examsubparts}
\end{examproblems}
\end{document}
